\section{Introduction}
\label{sec:intro}

Diffusion MRI turns each brain into a weighted graph: regions are nodes and white-matter tracts,
quantified by tractography-derived weights, are edges. PTSD has been linked to altered large-scale organisation
of such networks \citep{suo2017smallworld,akiki2017network}, making graph learning a natural tool
for studying it. Explanations of predictions may be useful for model auditing, yet message-passing
GNNs \citep{kipf2017gcn,velickovic2018gat} commonly require post-hoc tools
\citep{ying2019gnnexplainer}. GNAN \citep{bechler2024gnan} is algebraically decomposable,
but lets nodes interact through their \emph{hop distance}, which is uninformative in dense
connectomes when nearly every pair is one or two hops apart. Our cohort has only 60 positive screens,
the small-sample regime typical of clinical neuroimaging. We therefore test the mechanism where it
can be verified (synthetic data) and treat the real-data results, positive or negative, as an
empirical study with correspondingly limited inferential scope.

\paragraph{Research questions.} Methodologically, can a one-hop GNAN be extended to dense,
continuously weighted graphs while retaining an exact fixed-graph logit decomposition? Empirically,
in a pilot cohort of 869 subjects (60 PCL-5 positive screens), how does the extension compare with
linear and GNN baselines for binary probable-PTSD classification, and do region-level explanations
replicate on held-out subjects?

\paragraph{Contributions.}
\textbf{(C1)} W-GNAN, with region-specific shape functions and a monotone edge function $g(w)$,
whose logit decomposes into node and pairwise-interaction terms and which reduces to a specified
mean-aggregating one-hop GNAN on binary graphs (\S\ref{sec:method}).
\textbf{(C2)} A synthetic validation in which topology is uninformative by construction
(\S\ref{sec:res-synth}).
\textbf{(C3)} A benchmark on shared frozen splits with bootstrap CIs and selected-nuisance checks
(\S\ref{sec:res-bench}--\ref{sec:res-sanity}).
\textbf{(C4)} A region-level interpretability audit: a descriptive in-sample/held-out comparison,
label-permutation tests of hypothesis-motivated networks, and intrinsic/post-hoc agreement
(\S\ref{sec:res-interp}).
